\section{The Bayesian defect and causal compatibility}\label{sec:defect}
\subsection{The actual conditional experiment at an incoming label}
All conditional expectations below include the observer's previous independent randomizations. They do not include future audit outcomes. Fix a phase and suppress its index. For an incoming label $s$ of positive probability write $u_s=\E[X\mid S=s]$. Affinity of $\ell$ and $U$ gives, for every report event $E$,
\[
 \Pp(Y\in E\mid S=s)=\int_E\ell(y,u_s)\lambda(dy),\qquad
 \E[X'1_{\{Y\in E\}}\mid S=s]=\int_EU(y,u_s)\lambda(dy).
\]
Consequently, if $W=F(Y,u_S)$,
\begin{equation}\label{eq:bayes-barycenter}
 \E[X'\mid S,Y]=W.
\end{equation}
This is a conditional Bayes identity, not an interchange of a nonlinear map and an expectation under the old history law. The current report reweights that law by its likelihood. If $J$ is the new label, freshness of its update randomization implies
$\E[X'\mid S,Y,J]=W$. The same formula follows first for bounded simple reports and then by conditional integration; it does not condition a continuous law on an event of positive singleton probability.

For positive-mass pairs $(s,j)$ put
\[
 c_{sj}=\E[W\mid S=s,J=j],\qquad v_j=\E[W\mid J=j].
\]
Define, with zero-mass summands omitted,
\begin{equation}\label{eq:defects}
 D^{\rm row}=\E\norm{W-c_{SJ}}^2,\qquad
 D^{\rm pool}=\E\norm{c_{SJ}-v_J}^2.
\end{equation}
The tower property gives $v_J=\E[X'\mid J]$. Orthogonal conditional projections yield
\begin{align}
 \E\norm W^2&=D^{\rm row}+D^{\rm pool}+\E\norm{v_J}^2,\label{eq:two-projections}\\
 \E\norm{X'}^2&=\E\Phi(X),\qquad
 \E\norm W^2=\E\Phi(u_S).\label{eq:potential-moments}
\end{align}
Subtracting proves~\eqref{eq:main-identity}. In particular the potential's Jensen term is exactly
\begin{equation}\label{eq:residual-identity}
 \E\Phi(X)-\E\Phi(u_S)=\E\norm{X'-\E[X'\mid S,Y]}^2.
\end{equation}
This also proves its nonnegativity without differentiability. It measures the old information that remains inaccessible even after the new report has arrived.

\subsection{Perspective curvature}
For $q>0$, the function $(w,q)\mapsto\norm w^2/q$ is convex, since
\[
 \frac{\norm w^2}{q}=\sup_{z\in\R^k}\{2\langle z,w\rangle-q\norm z^2\}.
\]
The closed perspective at $(0,0)$ has value zero. The bounded prediction domain implies $U=0$ whenever $\ell=0$, so this convention agrees with the raw experiment. Integration and affine pullback prove convexity of $\Phi$ on $K$. For positive likelihood and a tangent direction $v$, direct differentiation gives
\begin{equation}\label{eq:perspective-hessian}
 \frac12D^2\!\left(\frac{\norm U^2}{\ell}\right)[v,v]
 =\frac{\norm{A v-F(c v)}^2}{\ell}
 =\ell\norm{D F[v]}^2.
\end{equation}
The assumed domination justifies integration and proves~\eqref{eq:hessian}. The formula is an acquired likelihood curvature: its weight is the physical report density, and its derivative is that of the normalized predictive update. It is not a separately specified Fisher metric.

When $\Phi$ is differentiable, define
\[
 D_\Phi(x,u)=\Phi(x)-\Phi(u)-D\Phi(u)[x-u].
\]
Conditioning on $S$ removes the last term and gives
\begin{equation}\label{eq:bregman}
 \E\Phi(X)-\E\Phi(u_S)=\E D_\Phi(X,u_S).
\end{equation}
If $\frac12D^2\Phi\le b I$ on every segment in $K$, Taylor's integral formula, including its factor $\int_0^1(1-s)ds=1/2$, implies
\begin{equation}\label{eq:jensen-upper}
 0\le \E D_\Phi(X,u_S)\le b\E\norm{X-u_S}^2.
\end{equation}
A measurable supporting affine function gives the analogous nonnegative identity on a nonsmooth domain whenever a subgradient with the required integrability is available. No such selection is needed for the unconditional identity~\eqref{eq:residual-identity}.

\subsection{Converse and the three compatibility questions}
For a finite measure $\zeta$ on a compact Euclidean set write
\[
 q_m(\zeta)=\inf_{a_1,\ldots,a_m}\int\min_j\norm{x-a_j}^2\,\zeta(dx).
\]
Centers may lie in the closed convex hull; compactness makes the infimum attainable. For any randomized assignment to at most $m$ labels and any assigned responses, conditional minimization gives distortion at least $q_m(\zeta)$. Conversely deterministic nearest-center assignment attains the center infimum. This argument applies to fractional cells as well as deterministic partitions.

At a fixed old label $s$, the law of $W$ is $\mu_{t,u_s}$, and the new label is a randomized assignment to at most $m_t$ values. Therefore
\[
 D^{\rm row}\ge\sum_s\Pp(S=s)q_{m_t}(\mu_{t,u_s}).
\]
Together with the nonnegative Jensen and pooling terms, this proves~\eqref{eq:main-converse}. The sharper statement retaining the Jensen term is immediate from the same equality.

For a \emph{fixed incoming law and a fixed one-step labeling}, equality in this sharper bound holds if and only if every positive-mass row attains its static quantization optimum and all row centroids using the same outgoing label agree. Necessity follows by summing nonnegative losses; sufficiency follows from~\eqref{eq:two-projections}. This separates static row optimality and zero pooling. It does not prove that such labelings can be chosen compatibly at every phase. Even an attained static optimum in every row need not admit a common outgoing labeling. The common-cover construction in the next section instead bounds the \emph{sum} of the two defects and proves a global matching rate without claiming local saturation.

In the state-independent linear case $\ell(y,u)=\ell(y)$, $F(y,u)=A(y)u$, and $\int A(y)^TA(y)\ell(y)\lambda(dy)=\kappa I$, one has $\Phi(u)=\kappa\norm u^2$. Thus~\eqref{eq:main-identity} reduces to
\[
 r_t=\kappa r_{t-1}+D^{\rm row}+D^{\rm pool},\qquad
 e_t=\kappa e_{t-1}-D^{\rm row}-D^{\rm pool}.
\]
The exact energy and circle results of Companion G are recovered with their original assumptions. Zero-energy rows contribute zero; a later zero product factor does not certify earlier saturation. Their directional coefficient is an infimum, and equality requires an actual minimizing direction and row encoder, not merely approximate attainment.

\begin{example}[An informative report cannot be treated as a neutral transport]\label{ex:binary}
Prepare $X\in\{-1,1\}$ with equal probabilities and retain one old label. Let $0<\eta<1$ and report $Y\in\{-1,1\}$ with $\Pp(Y=X\mid X)=(1+\eta)/2$. A two-label update stores $Y$. Its retained posterior is $\eta Y$, so its retained energy rises from zero to $\eta^2$. For a prior mean $u$, the report likelihood is $(1+\eta yu)/2$ and the posterior mean is $(u+\eta y)/(1+\eta yu)$. Its potential is
\[
 \Phi(u)=\frac12\sum_{y=\pm1}\frac{(u+\eta y)^2}{1+\eta yu}.
\]
Here $\Phi(0)=\eta^2$, whereas $\Phi(1)=\Phi(-1)=1$. The identity gives residual error $1-\eta^2$ and zero row/pooling loss. It explains the increase in information through the likelihood term; a state-independent energy recursion would miss it.
\end{example}
